\begin{textsample}{NEG Dim 3 – human – Score 1.00 – sp/sp\_inf\_e\_ef\_intro\_6\_p8.txt}  \label{ex:f3_neg_003}
Não é demais reforçar que as práticas de ensino e de aprendizagem que consideram o estudante em sua integralidade estão longe de práticas que normatizam comportamentos, rotulam ou buscam adequar os estudantes a um modelo ideal de pessoa.
A Educação Integral, como funda¬mento pedagógico, demonstra o interesse do Currículo Paulista em atender às necessidades de ensino e de aprendizagem pelo olhar sistêmico — por parte dos profissionais da educação — para essas aprendizagens e o modo como elas se apresentam em nossa sociedade.

% matched lemmas: 
\end{textsample}
